\begin{document}
\vspace*{0.2in}

\begin{flushleft}
{\Large
\textbf{Checking agent-run interpretability of a single-cell foundation model: a controlled evaluation of method specifications and a reviewer-derived audit checklist, with a MaxToki-217M case study}
}
\newline
\\
Ihor Kendiukhov\textsuperscript{1*}
\\
\bigskip
\textbf{1} Institute of Medical Genetics and Applied Genomics, University of T\"ubingen, T\"ubingen, Germany
\\
\bigskip

* ihor.kendiukhov@student.uni-tuebingen.de

\end{flushleft}

\section*{Abstract}
We evaluated one way to check interpretability analyses that
large-language-model agents run on biological foundation models. Each
method was written as an agent specification; outputs were audited with
a ten-item checklist of peer reviewers' objections to earlier studies.
An agent used both to run eight interpretability pipelines on the
single-cell foundation model MaxToki-217M and to audit them. We confirmed 149 errors in the deployment's code and outputs, 13
of them critical, including intervention code that deleted a transformer
block and model inputs encoded on the wrong scale. The agent's own audit report found 10\% of a fixed set of 68
errors that predated it, and no critical one; fresh agents in a blind re-audit found 41--45\%, with or
without the checklist. In a controlled study of four known-answer tasks,
all 40 Claude Opus~5.5 runs reached the correct conclusion, with or
without the specification. Executors handled 94.5\% of pitfalls; review and repair added 2.0--2.8 points; generic reviews
turned 3 of 40 correct conclusions wrong. In 11 held-out analyses from
other projects, each in a flawed and a corrected version, every review
of a flawed version found its error, but about 80\% of
reviews of corrected versions claimed a serious error that was not one.
With Opus~5.5, no study showed the checklist finding more errors or
raising fewer false alarms than an equal-budget generic review. Of the 13 critical errors, author-directed agent checks found 8 first and
blind re-audits 3. After correction, MaxToki-217M's attention scores,
sparse-autoencoder features and traced circuits did not beat gene-level
baselines; attention matched co-expression and beat rewired TRRUST
networks in RPE1 cells, not K562. Its gene embeddings aligned with
scGPT's and more with Geneformer's. Hidden states showed no
developmental order beyond cell type; 14 of 15 features steered a maturity score no more than random features.

\clearpage
\newgeometry{top=0.85in,left=1in,right=1in,footskip=0.75in}
\fancyheadoffset[L]{0pt}
\fancyfootoffset[L]{0pt}
\linenumbers
